\section{Limitations and Claim Boundary}
\label{sec:ieee-limitations}

The ORIUS argument is strongest when the claim boundary remains explicit. This draft
does not claim universal optimality, equal-theorem-depth parity, universal hardware
closure, or flat deployment readiness across all supported domains. It claims one
universal runtime safety layer, one reusable action-semantics grammar for degraded
observation, and one evidence ladder that determines how much each domain can currently
support.

That distinction matters because the architecture already travels farther than the
tracked evidence. Battery is the primary validation domain. Autonomous vehicles and
healthcare monitoring are bounded validation domains under explicit contracts. Future
additional domains remain outside the current evidence boundary. The manuscript
therefore uses the active three-domain program to show the reach of the runtime
contract while preserving the boundary between artifact-supported results and broader
architectural framing.

The bounded-universal position is still scientifically valuable. It makes ORIUS a real
cross-domain safety layer today, while giving future closure work a precise target:
new domains must supply defended replay, runtime, governance, and post-repair
improvement surfaces on the governing safety object. Those are evidence obligations,
not conceptual holes in the runtime grammar itself. The adversarial OQE mode
partially addresses the reliability-spoofing limitation, but full game-theoretic
completeness across all attack strategies remains an open evidence obligation.

\IfFileExists{reports/publication/security_governance_ablation_matrix.tex}{%
  \begin{table*}[t]
\centering
\caption{Governance and security ablations for the ORIUS runtime-release contract. These are fail-closed ablations, not performance benchmarks: each row removes or corrupts a release-governance precondition and records the expected denial path.}
\label{tab:security-governance-ablations}
\footnotesize
\renewcommand{\arraystretch}{1.10}
\begin{tabularx}{\textwidth}{p{0.14\textwidth} p{0.20\textwidth} p{0.18\textwidth} p{0.18\textwidth} X}
\toprule
\textbf{Ablation} & \textbf{Removed or corrupted component} & \textbf{Expected runtime response} & \textbf{Evidence surface} & \textbf{Interpretation} \\
\midrule
Missing model hash & Artifact hash manifest absent in strict mode & Refuse load before deserialization & model-hash test; deployment-security validator & Strict release evidence cannot depend on unverifiable model binaries. \\
Bad certificate signature & Certificate payload or signature tampered & Verification fails closed & certificate-full test; schema validator & Certificate-backed release is meaningful only when invalid signatures deny release. \\
Stale device timestamp & Device request timestamp outside allowed skew & Reject telemetry or acknowledgement & IoT HMAC test; deployment-security validator & Device identity is enforced as a runtime precondition. \\
Replayed device nonce & Previously accepted device nonce reused & Reject replayed device request & IoT HMAC test; deployment-security validator & Strict IoT mode includes replay protection. \\
Invalid fallback action & Fallback fails the domain postcondition & Release-contract validation fails & release-contract test and validator & Fallback remains a typed safety obligation, not an unconditional escape. \\
Corrupted release witness & Compact evidence omits canonical fields & Release-contract validator fails & release-witness corruption test and validator & Claim-facing evidence must expose common release fields across promoted domains. \\
\bottomrule
\end{tabularx}
\end{table*}
%
}{}

\IfFileExists{reports/publication/orius_failure_modes_falsification_table.tex}{%
  \begin{table*}[t]
\centering
\caption{Falsification boundaries for the ORIUS claim. These rows define what would break the present bounded-predeployment claim and how the repo currently guards or scopes that risk.}
\label{tab:orius-failure-falsification}
\footnotesize
\renewcommand{\arraystretch}{1.10}
\begin{tabularx}{\textwidth}{p{0.18\textwidth} p{0.29\textwidth} p{0.24\textwidth} X}
\toprule
\textbf{Failure mode} & \textbf{What would falsify the current claim} & \textbf{Current guardrail} & \textbf{Boundary or next step} \\
\midrule
Coverage miss exceeds $\alpha$ & Empirical or certified coverage miss rate exceeds the theorem/reporting budget on a promoted row & Coverage/calibration artifacts plus the T11 coverage obligation & Recalibrate uncertainty sets or demote the affected row until coverage is restored. \\
Safe-set adapter is wrong & Domain audit shows $C(x)$ excludes a truly safe action or includes an unsafe action & Adapter-validity assumption and domain runtime-contract witness & Correct the adapter and rerun domain-specific runtime evidence. \\
Fallback is not admissible & Fallback violates true-state safety or operational admissibility under a promoted scenario & Release-contract postcondition tests and fallback-validity theorem rows & Replace fallback, deny release, or narrow the claim boundary. \\
Repair leaves the safe core & Repaired action fails membership in the uncertainty-set safe-action intersection & Canonical release evidence and T11 repair-membership obligation & Fail closed and rerun repair validation before promotion. \\
Audit chain is broken & Certificate hash, signature, or append-only chain cannot be verified for a promoted release row & Certificate schema and deployment-security validators & Do not use the row as claim-governing evidence until provenance is repaired. \\
Utility collapses to shutdown & ORIUS matches a fail-safe reference only by always braking, alerting, or shutting down & Utility-preserving safety scorecard & Report the row as safe but vacuous; improve controller/fallback policy before claiming non-vacuity. \\
Closed-loop dynamics diverge & nuPlan, CARLA, or physical loop reveals behavior not represented by bounded replay/surrogate dynamics & Validation-boundary table and external-validation manifest & Promote only after closed-loop evidence exists; otherwise keep the replay boundary explicit. \\
Clinical context differs & Prospective/site clinical use has different workflow, population, or utility costs than retrospective monitoring & Clinical boundary table and source/time-holdout evidence & Require prospective/site validation before clinical-use claims. \\
\bottomrule
\end{tabularx}
\end{table*}
%
}{}
